\documentclass[aspectratio=169]{beamer}

\usetheme{Madrid}
\definecolor{unsblue}{RGB}{30,60,120}
\definecolor{unslightblue}{RGB}{220,230,245}
\definecolor{alertred}{RGB}{180,30,30}
\usecolortheme[named=unsblue]{structure}
\setbeamercolor{palette primary}{bg=unsblue, fg=white}
\setbeamercolor{palette secondary}{bg=unsblue!80, fg=white}
\setbeamercolor{palette tertiary}{bg=unsblue!60, fg=white}
\setbeamercolor{palette quaternary}{bg=unsblue, fg=white}
\setbeamercolor{block title}{bg=unsblue, fg=white}
\setbeamercolor{block body}{bg=unslightblue}
\setbeamercolor{block title alerted}{bg=alertred, fg=white}
\setbeamercolor{block body alerted}{bg=red!8}
\setbeamercolor{block title example}{bg=unsblue!60!black, fg=white}
\setbeamercolor{block body example}{bg=unsblue!8}

\usepackage[utf8]{inputenc}
\usepackage[T1]{fontenc}
\usepackage{babel}
\babelprovide[main, import]{spanish}
\usepackage{amsmath, amssymb}
\usepackage{bm}
\usepackage{booktabs}
\usepackage{multirow}

\newcommand{\R}{\mathbb{R}}
\newcommand{\bx}{\bm{x}}

\title[MAD1 -- Unidad 4]{Aprendizaje Supervisado}
\subtitle{Clase 23: Redes Neuronales, Evaluación y Pipelines}
\author[J. Bavio]{Dr.\ José Bavio}
\institute[UNS]{Departamento de Matemática\\Universidad Nacional del Sur}
\date{2026}

\begin{document}

\begin{frame}
  \titlepage
\end{frame}

\begin{frame}{Contenido de la clase}
  \tableofcontents
\end{frame}

% ═════════════════════════════════════════════════════════════════════════════
\section{Redes Neuronales Artificiales}
% ═════════════════════════════════════════════════════════════════════════════

\begin{frame}{La neurona artificial}
  \begin{block}{Idea}
    Una neurona recibe $p$ entradas, las combina linealmente y aplica una función no lineal:
    \[
      a = \sigma\!\left(\sum_{j=1}^{p} w_j x_j + b\right) = \sigma(\bm{w}^\top \bx + b)
    \]
  \end{block}

  \medskip
  \begin{exampleblock}{Analogía biológica (solo para intuición)}
    En el cerebro, una neurona ``se activa'' si la suma de señales que recibe supera un umbral.
    La función de activación $\sigma$ cumple el mismo rol: decide si la neurona ``se prende'' o no.\\[4pt]
    La analogía no es perfecta, pero ayuda a recordar la estructura.
  \end{exampleblock}
\end{frame}

\begin{frame}{Perceptrón Multicapa (MLP)}
  \begin{block}{Definición}
    Un \textbf{MLP} (\textit{Multilayer Perceptron}) apila capas de neuronas:
    \begin{itemize}
      \item \textbf{Capa de entrada:} dimensión $p$ (una neurona por feature).
      \item \textbf{Capas ocultas:} transformaciones no lineales sucesivas.
      \item \textbf{Capa de salida:} produce la predicción final.
    \end{itemize}
  \end{block}

  \medskip
  Salida de la capa $l$:
  \[
    \bm{h}^{(l)} = \sigma^{(l)}\!\left(\bm{W}^{(l)} \bm{h}^{(l-1)} + \bm{b}^{(l)}\right)
  \]

  \begin{exampleblock}{¿Por qué capas ocultas?}
    Cada capa aprende representaciones más abstractas de los datos.
    La primera capa puede detectar bordes; la segunda, formas; la tercera, objetos.
    Con suficientes capas y neuronas, un MLP puede aproximar cualquier función continua
    (\textit{teorema de aproximación universal}).
  \end{exampleblock}
\end{frame}

\begin{frame}{Funciones de activación}
  \begin{center}
  \begin{tabular}{llll}
    \toprule
    \textbf{Nombre} & \textbf{Fórmula} & \textbf{Rango} & \textbf{Uso} \\
    \midrule
    Sigmoide    & $1/(1+e^{-z})$          & $(0,1)$    & Salida binaria \\
    Tanh        & $\tanh(z)$              & $(-1,1)$   & Capas ocultas \\
    ReLU        & $\max(0,z)$             & $[0,\infty)$ & Estándar hoy \\
    Leaky ReLU  & $\max(\alpha z,z)$      & $\R$       & Evita neuronas muertas \\
    Softmax     & $e^{z_k}/\sum_j e^{z_j}$ & $(0,1)^K$ & Salida multiclase \\
    \bottomrule
  \end{tabular}
  \end{center}

  \medskip
  \begin{exampleblock}{¿Por qué ReLU domina hoy?}
    La sigmoide y la tanh ``aplastan'' los gradientes en sus extremos (gradiente $\approx 0$):
    el entrenamiento se detiene. ReLU no tiene ese problema.
    Además es muy simple de calcular: $\max(0,z)$.
  \end{exampleblock}
\end{frame}

\begin{frame}{Entrenamiento: descenso por gradiente}
  Los parámetros $\bm{\theta} = \{\bm{W}^{(l)}, \bm{b}^{(l)}\}$ se ajustan minimizando la función de pérdida.

  \medskip
  \textbf{Descenso por gradiente estocástico (SGD):}
  \[
    \bm{\theta} \leftarrow \bm{\theta} - \eta\, \nabla_{\bm{\theta}}\, \mathcal{L}
  \]
  donde $\eta$ es la \textbf{tasa de aprendizaje} (\textit{learning rate}).

  \medskip
  \begin{exampleblock}{Intuición}
    Imaginar que estás en una montaña con niebla y querés bajar al valle.
    El gradiente te dice la dirección de máxima subida; ir en la dirección opuesta
    (descenso por gradiente) te lleva hacia abajo.
    El learning rate es cuán grande es cada paso que das.
  \end{exampleblock}

  \medskip
  \begin{alertblock}{Cuidado con $\eta$}
    $\eta$ muy grande: el modelo ``salta'' y no converge.\\
    $\eta$ muy pequeño: convergencia lentísima.
  \end{alertblock}
\end{frame}

\begin{frame}{Backpropagation}
  \begin{block}{¿Cómo calcular $\nabla_{\bm{\theta}}\,\mathcal{L}$?}
    \textbf{Backpropagation} aplica la regla de la cadena capa por capa,
    de la salida hacia la entrada, calculando eficientemente el gradiente
    respecto a cada parámetro de la red.
  \end{block}

  \medskip
  \begin{exampleblock}{Aclaración importante}
    Backprop no es el algoritmo de optimización: es el método para calcular el gradiente.
    El algoritmo de optimización (SGD, Adam, etc.) usa ese gradiente.\\[4pt]
    Es como la diferencia entre calcular la pendiente de una montaña (backprop)
    y la estrategia de descenso que usás (SGD, Adam).
  \end{exampleblock}

  \medskip
  \textbf{Adam} (Kingma \& Ba, 2015) es el optimizador más usado hoy:
  combina momentum con adaptación individual de la tasa por parámetro.
\end{frame}

\begin{frame}{Regularización en redes neuronales}
  \begin{columns}
    \column{0.5\textwidth}
    \begin{block}{Dropout}
      Durante el entrenamiento, cada neurona se desactiva con probabilidad $p_{\text{drop}}$.
      Fuerza redundancia: la red no puede depender de ninguna neurona individual.\\[4pt]
      \textit{Como estudiar para el examen sin saber de antemano qué páginas del libro van a estar disponibles.}
    \end{block}

    \column{0.5\textwidth}
    \begin{block}{Early stopping}
      Entrenar hasta que el error de validación deja de mejorar.
      Guardar los pesos del mejor epoch.\\[4pt]
      \textit{Parar de estudiar cuando ya no aprendés más, antes de que empieces a ``aprender'' los ruidos.}
    \end{block}
  \end{columns}

  \medskip
  \begin{block}{Otras técnicas}
    \begin{itemize}
      \item Regularización $L_2$ (\textit{weight decay}): penaliza pesos grandes.
      \item Batch normalization: normaliza activaciones, acelera el entrenamiento.
    \end{itemize}
  \end{block}
\end{frame}

% ═════════════════════════════════════════════════════════════════════════════
\section{Métricas de Evaluación}
% ═════════════════════════════════════════════════════════════════════════════

\begin{frame}{Matriz de confusión}
  Para clasificación binaria (positivo / negativo):

  \medskip
  \begin{center}
  \begin{tabular}{cc|cc}
    & & \multicolumn{2}{c}{\textbf{Predicho}} \\
    & & Positivo & Negativo \\
    \hline
    \multirow{2}{*}{\textbf{Real}} & Positivo & VP & FN \\
    & Negativo & FP & VN \\
  \end{tabular}
  \end{center}

  \medskip
  \begin{itemize}
    \item \textbf{VP} (verdadero positivo): predijo positivo y era positivo. Bien.
    \item \textbf{VN} (verdadero negativo): predijo negativo y era negativo. Bien.
    \item \textbf{FP} (falso positivo): predijo positivo pero era negativo. \textit{Error tipo I.}
    \item \textbf{FN} (falso negativo): predijo negativo pero era positivo. \textit{Error tipo II.}
  \end{itemize}
\end{frame}

\begin{frame}{Métricas derivadas}
  \[
    \text{Accuracy} = \frac{\text{VP}+\text{VN}}{n}
    \qquad
    \text{Precisión} = \frac{\text{VP}}{\text{VP}+\text{FP}}
    \qquad
    \text{Recall} = \frac{\text{VP}}{\text{VP}+\text{FN}}
  \]
  \[
    F_1 = 2\cdot\frac{\text{Precisión}\times\text{Recall}}{\text{Precisión}+\text{Recall}}
  \]

  \medskip
  \begin{exampleblock}{Intuición de Precisión y Recall}
    \textbf{Precisión:} ``De todos los que el modelo acusó de positivos, ¿cuántos realmente lo eran?''\\
    \textbf{Recall:} ``De todos los positivos reales, ¿cuántos detectó el modelo?''\\[4pt]
    En diagnóstico médico: recall alto es crítico (no quiero perder enfermos).\\
    En correo spam: precisión alta importa más (no quiero tirar mail legítimo).
  \end{exampleblock}
\end{frame}

\begin{frame}{El problema de la accuracy con clases desbalanceadas}
  \begin{alertblock}{Ejemplo: detección de fraude}
    99\% de las transacciones son legítimas, 1\% son fraude.\\
    Un clasificador que predice SIEMPRE ``no fraude'' tiene accuracy = 99\%.\\
    Sin embargo, \textbf{no detecta ningún fraude}.
  \end{alertblock}

  \medskip
  \begin{block}{Soluciones}
    Con clases desbalanceadas, preferir:
    \begin{itemize}
      \item $F_1$ o $F_\beta$ (con $\beta > 1$ si recall importa más)
      \item AUC-ROC (Area Under the ROC Curve)
      \item AUC-PR (Area Under the Precision-Recall Curve)
    \end{itemize}
  \end{block}
\end{frame}

\begin{frame}{Curva ROC y AUC}
  La \textbf{curva ROC} grafica:
  \begin{itemize}
    \item Eje $y$: TPR (Recall, Sensibilidad) $= \text{VP}/(\text{VP}+\text{FN})$
    \item Eje $x$: FPR $= \text{FP}/(\text{FP}+\text{VN})$
  \end{itemize}
  para todos los posibles umbrales $\tau \in [0,1]$.

  \medskip
  \begin{block}{AUC-ROC}
    \begin{itemize}
      \item AUC $= 1$: clasificador perfecto.
      \item AUC $= 0.5$: clasificador aleatorio (diagonal).
      \item AUC $< 0.5$: peor que aleatorio.
    \end{itemize}
  \end{block}

  \medskip
  \begin{exampleblock}{Interpretación probabilística}
    AUC = probabilidad de que el modelo asigne un score más alto a un positivo aleatorio
    que a un negativo aleatorio.\\
    No depende del umbral de decisión: mide la capacidad discriminante global.
  \end{exampleblock}
\end{frame}

% ═════════════════════════════════════════════════════════════════════════════
\section{Clases Desbalanceadas y SMOTE}
% ═════════════════════════════════════════════════════════════════════════════

\begin{frame}{¿Por qué el desbalance es un problema?}
  \begin{block}{Efecto sobre el entrenamiento}
    Si el 95\% de los datos son clase mayoritaria, el modelo aprende rápido a
    ignorar la clase minoritaria: el error en la mayoría domina la función de pérdida.
    El modelo ``toma el atajo'' de predecir siempre la clase mayoritaria.
  \end{block}

  \medskip
  \textbf{Estrategias principales:}
  \begin{enumerate}
    \item Resampling (oversampling / undersampling)
    \item Ajuste de pesos de clase (\texttt{class\_weight})
  \end{enumerate}
\end{frame}

\begin{frame}{SMOTE}
  \textbf{SMOTE} (\textit{Synthetic Minority Oversampling Technique}, Chawla et al., 2002)
  genera observaciones \textbf{sintéticas} de la clase minoritaria:

  \medskip
  \begin{enumerate}
    \item Para cada punto $\bx_i$ de la clase minoritaria, encontrar sus $k$ vecinos más cercanos.
    \item Generar un punto sintético interpolando:
    \[
      \tilde{\bx} = \bx_i + \lambda\,(\bx_{nn} - \bx_i), \quad \lambda \sim \mathcal{U}(0,1)
    \]
  \end{enumerate}

  \medskip
  \begin{exampleblock}{Coloquialmente}
    En lugar de copiar y pegar los mismos ejemplos minoritarios (oversampling simple),
    SMOTE ``inventa'' ejemplos nuevos que están entre dos ejemplos reales de la clase.
    Añade variedad sin salir de la región donde ya existen ejemplos.
  \end{exampleblock}

  \medskip
  \begin{alertblock}{Regla fundamental}
    SMOTE se aplica \textbf{solo sobre el conjunto de entrenamiento}.
    Nunca antes de dividir los datos: eso produce data leakage.
  \end{alertblock}
\end{frame}

% ═════════════════════════════════════════════════════════════════════════════
\section{Pipelines y Grid Search}
% ═════════════════════════════════════════════════════════════════════════════

\begin{frame}{¿Qué es un pipeline?}
  \begin{block}{Definición}
    Un \textbf{pipeline} encadena pasos de preprocesamiento y modelado en un único objeto.
    Cada paso transforma los datos y los pasa al siguiente.
  \end{block}

  \medskip
  Ejemplo típico:
  \[
    \text{Datos} \xrightarrow{\text{imputar}} \xrightarrow{\text{estandarizar}} \xrightarrow{\text{SMOTE}} \xrightarrow{\text{modelo}} \text{Predicción}
  \]

  \medskip
  \begin{exampleblock}{¿Por qué es importante?}
    Sin pipeline, es fácil cometer el error de:
    \begin{enumerate}
      \item Estandarizar todos los datos.
      \item \textit{Luego} dividir en train/test.
    \end{enumerate}
    Resultado: el test set ``contaminó'' los parámetros de la estandarización.
    Eso es \textbf{data leakage}: el modelo vio información del test durante el entrenamiento.
  \end{exampleblock}
\end{frame}

\begin{frame}{Data Leakage}
  \begin{alertblock}{El error más silencioso del ML}
    El \textbf{data leakage} ocurre cuando información del futuro o del test set
    contamina el entrenamiento del modelo.\\
    El modelo parece funcionar muy bien en evaluación, pero falla en producción.
  \end{alertblock}

  \medskip
  \begin{block}{Formas comunes de leakage}
    \begin{itemize}
      \item Normalizar/estandarizar con estadísticos calculados sobre todo el dataset.
      \item Aplicar SMOTE antes de dividir train/test.
      \item Seleccionar features con toda la información antes de dividir.
      \item Usar features que ``vienen del futuro'' (ej.: resultado de un evento para predecirlo).
    \end{itemize}
  \end{block}

  \medskip
  \begin{exampleblock}{Solución}
    Usar pipelines con validación cruzada: cada transformación se ajusta solo con los datos de entrenamiento de cada fold.
  \end{exampleblock}
\end{frame}

\begin{frame}{Grid Search y búsqueda de hiperparámetros}
  \begin{block}{Grid Search}
    Define una grilla de valores para cada hiperparámetro y evalúa todas las combinaciones
    usando validación cruzada sobre el conjunto de entrenamiento.
  \end{block}

  \medskip
  \begin{center}
  \begin{tabular}{ll}
    \toprule
    \textbf{Método} & \textbf{Descripción} \\
    \midrule
    Grid Search   & Todas las combinaciones (exhaustivo, costoso) \\
    Random Search & Muestrea combinaciones al azar (más eficiente) \\
    Bayesiano     & Guía la búsqueda con un modelo probabilístico \\
    \bottomrule
  \end{tabular}
  \end{center}

  \medskip
  \begin{alertblock}{Regla de oro}
    La evaluación con grid search usa el \textbf{conjunto de entrenamiento} (con CV interna).
    El conjunto de prueba queda reservado para la evaluación final.
    Mirar el test set durante la búsqueda de hiperparámetros invalida su uso como estimador imparcial.
  \end{alertblock}
\end{frame}

\begin{frame}{Pipeline + Grid Search: la combinación correcta}
  Con \texttt{GridSearchCV} en scikit-learn sobre un pipeline:

  \medskip
  \begin{enumerate}
    \item En cada fold de CV:
    \begin{itemize}
      \item Los pasos del pipeline se \textbf{ajustan} sobre el fold de entrenamiento.
      \item Se \textbf{aplican} sobre el fold de validación.
    \end{itemize}
    \item Se elige el conjunto de hiperparámetros con mejor CV.
    \item Se re-entrena el pipeline completo con todos los datos de entrenamiento.
    \item Se evalúa \textbf{una sola vez} sobre el test set.
  \end{enumerate}

  \medskip
  \begin{exampleblock}{Coloquialmente}
    Es el equivalente a preparar el examen con ejercicios de práctica
    sin mirar las respuestas del examen real hasta el día del examen.
    Una sola vez. Sin volver a estudiar después de ver el resultado.
  \end{exampleblock}
\end{frame}

% ─────────────────────────────────────────────────────────────────────────────
\begin{frame}{Resumen de la Unidad 4}
  \begin{center}
  \begin{tabular}{lll}
    \toprule
    \textbf{Modelo} & \textbf{Fortaleza} & \textbf{Limitación clave} \\
    \midrule
    Árbol de decisión & Interpretable      & Alta varianza \\
    Random Forest     & Robusto, OOB       & Caja negra \\
    SVM               & Alta dimensión     & Costoso en $n$ grande \\
    KNN               & Sin entrenamiento  & Maldición dimensionalidad \\
    MLP               & Muy flexible       & Requiere muchos datos \\
    \bottomrule
  \end{tabular}
  \end{center}

  \bigskip
  \textbf{Temas transversales:}
  \begin{itemize}
    \item Bias-variance tradeoff y validación cruzada (siempre).
    \item Métricas apropiadas al problema (accuracy no siempre sirve).
    \item Pipelines para evitar data leakage.
    \item SMOTE para clases desbalanceadas.
  \end{itemize}
\end{frame}

\end{document}
